% !PRINTED 3
% Source chapter opening suppresses the running head and folio. Logical printed page 3 is
% fixed by the following explicit folio 4 and the preceding prelim sequence. The "1*" at
% the foot is a printer's signature, not a folio or body text, and is intentionally omitted.
\weberpage{3}

\begin{center}
\gesperrt{First Chapter.}
\end{center}

\begin{center}
\textbf{The Elliptic Integrals.}
\end{center}

\begin{center}
\rule{3cm}{0.4pt}
\end{center}

\begin{center}
\textbf{\S\ 1. \gesperrt{Definition of the Elliptic Integrals.}}
\end{center}

When the systematic exposition of the integral calculus has advanced to the point at which algebraic integrals containing the square root of a function of the first or second degree are reduced to integrals of rational functions, the learner meets at this place a barrier which he is not in a position to pass with the means that stand at his command up to that point. The striving after an extension of those means, so as to make the next class of integrals accessible to research as well, was historically the occasion of a more thorough study of the elliptic integrals and of the introduction of the elliptic functions; and it is also the most natural and most intelligible starting point for one who is to be introduced to the theory of these functions for the first time. It shall therefore be our next task to make ourselves acquainted with the elliptic integrals and with their most important properties.

The definition of an elliptic integral from which we wish to start is the following: let $f(x)$ denote an integral rational function of the third or fourth degree with four distinct roots
\weq{1}{f(x) = a_0 x^4 + 4 a_1 x^3 + 6 a_2 x^2 + 4 a_3 x + a_4,}
in which $a_0$ and $a_1$ do not both vanish at once; let $\Phi(x, y)$ further be an arbitrary integral or fractional rational function of the two arguments $x, y$. Then
\[
\int \Phi(x, \sqrt{f(x)})\, dx
\]
is the general \gesperrt{elliptic integral}, and
\[
\Phi(x, \sqrt{f(x)})\, dx
\]
%CONT>>
